\chapter[Marco Legal]{Marco Legal}
\label{ch:marco-legal}

\insertminitoc
\parindent0pt

El desarrollo de una aplicación de reconocimiento de lengua de señas para niños sordos se
enmarca en un conjunto de disposiciones jurídicas que protegen los derechos de las personas
con discapacidad y promueven su inclusión educativa y comunicativa. Este capítulo presenta
los instrumentos legales internacionales y nacionales que fundamentan y legitiman el
presente proyecto, así como los estándares técnicos de accesibilidad que orientan el diseño
del sistema.

\section{Derecho a la Comunicación y Educación Inclusiva en el Marco Internacional}

La Convención sobre los Derechos de las Personas con Discapacidad (\acrshort{cdpd}), adoptada por la
Asamblea General de las Naciones Unidas en 2006 y ratificada por Honduras, establece en su
artículo 9 la obligación de los Estados de asegurar el acceso de las personas con
discapacidad a la información y las comunicaciones, incluidas las tecnologías de la
información. El artículo 21 reconoce el derecho a la libertad de expresión mediante las
lenguas de señas y otros modos y formatos de comunicación accesibles, mientras que el
artículo 24 exige que los Estados garanticen una educación inclusiva a todos los niveles,
facilitando el aprendizaje de la lengua de señas y promoviendo la identidad lingüística de
las personas sordas \cite{onu_cdpd_2006}.

La Agenda 2030 para el Desarrollo Sostenible, aprobada por las Naciones Unidas en 2015,
incorpora la discapacidad como eje transversal en sus metas. El Objetivo de Desarrollo
Sostenible 4 (Educación de Calidad) establece el compromiso de garantizar una educación
inclusiva, equitativa y de calidad para todas las personas. El Objetivo 10 (Reducción de
las Desigualdades) promueve la inclusión social y económica de las personas con
discapacidad independientemente de su condición \cite{onu_agenda2030_2015}. El presente
proyecto contribuye directamente a ambos objetivos al proporcionar una herramienta
tecnológica que facilita la comunicación entre niños sordos y personas oyentes en contextos
educativos.

La Declaración de Salamanca, adoptada por la UNESCO en 1994, constituye un referente
fundamental en materia de educación inclusiva. Este instrumento proclama que los sistemas
educativos deben diseñarse para atender la diversidad de necesidades de aprendizaje, y que
los niños con necesidades educativas especiales deben tener acceso a las escuelas ordinarias
como medio más eficaz para combatir actitudes discriminatorias y construir una sociedad
inclusiva \cite{unesco_salamanca_1994}.

\section{Legislación Hondureña en Materia de Discapacidad y Lengua de Señas}

La Constitución de la República de Honduras, en su artículo 60, declara punible toda
discriminación. El artículo 151 establece que la educación es función esencial del Estado,
y el artículo 169 dispone que el Estado sostendrá y fomentará la educación de las personas
con discapacidad \cite{constitucion_honduras_1982}. Estas disposiciones constitucionales
proporcionan el fundamento jurídico superior sobre el cual se apoyan las leyes específicas
en materia de discapacidad.

La Ley de Equidad y Desarrollo Integral para las Personas con Discapacidad (Decreto
160-2005) tiene como finalidad garantizar el disfrute pleno de los derechos de las personas
con discapacidad y promover su desarrollo integral dentro de la sociedad. Esta ley prohíbe
toda forma de discriminación directa o indirecta, establece principios de autodeterminación
y accesibilidad universal, y obliga a los órganos rectores de educación a asegurar el
acceso efectivo de las personas con discapacidad al sistema educativo \cite{decreto_160_2005}.
En su artículo 12, la ley establece que los medios de comunicación deben facilitar el
acceso a la información para las personas con discapacidad, lo que respalda el desarrollo
de herramientas tecnológicas como la que propone este proyecto.

La Ley de la Lengua de Señas Hondureña (Decreto 321-2013) reconoce al \acrshort{lesho} como medio de
comunicación oficial en el territorio nacional para las personas sordas, con discapacidad
auditiva y sordociegas, y establece la obligación del Estado de
promover su uso, enseñanza y difusión en todos los ámbitos sociales y educativos
\cite{congreso_decreto_2013}. Esta ley es de particular relevancia para el presente
proyecto, ya que legitima el desarrollo de tecnologías orientadas al reconocimiento y la
difusión de \acrshort{lesho} como instrumento de comunicación y de acceso a la educación para la
comunidad sorda hondureña.

\section{Normativa Técnica de Accesibilidad Aplicable al Sistema}

Las Pautas de Accesibilidad para el Contenido Web (\acrshort{wcag}) 2.1, publicadas por el World Wide
Web Consortium (\acrshort{w3c}) en 2018, constituyen el estándar internacional de referencia para el
diseño de contenido digital accesible. Aunque fueron concebidas para contenido web, sus
cuatro principios (perceptible, operable, comprensible y robusto) son aplicables al diseño
de aplicaciones móviles \cite{w3c_wcag_2018}. Para el presente sistema, el principio de
perceptibilidad resulta determinante: toda retroalimentación del sistema debe ser visual o
háptica, dado que el usuario principal es un niño sordo. Los criterios de contraste de
color, tamaño de elementos interactivos y alternativas textuales para contenido no textual
orientan directamente las decisiones de diseño de la interfaz de la aplicación.

\section{Protección de los Datos de las Personas Participantes}

El sistema propuesto capta imágenes de personas mediante la cámara del dispositivo, y su
construcción exigió grabar a colaboradores ejecutando señas. Ambas circunstancias sitúan
al proyecto en el terreno de la protección de datos personales, por lo que corresponde
precisar el marco jurídico aplicable en Honduras.

A la fecha de este trabajo, Honduras no cuenta con una ley integral de protección de datos
personales en vigor. Existe un anteproyecto que ha sido presentado ante el Congreso
Nacional, pero no ha completado su trámite legislativo. En ausencia de esa norma
específica, la protección se sustenta en el texto constitucional.

La Constitución de la República de Honduras, en su artículo 76, garantiza el derecho al
honor, a la intimidad personal y familiar y a la propia imagen. El artículo 100 establece
la inviolabilidad y el secreto de las comunicaciones. Y el artículo 182, en su numeral
segundo, reconoce la garantía de hábeas data, que otorga a toda persona un recurso para
conocer y controlar la información que sobre ella conste en registros o bancos de datos
\cite{constitucion_honduras_1982}. El derecho a la propia imagen es el que incide de forma
más directa en este proyecto, tanto por la captación de video durante la construcción del
conjunto de datos como por el uso de la cámara en la aplicación.

Frente a ese marco, el diseño del sistema adopta dos garantías técnicas que van más allá
de lo que exigiría una norma. La primera es que los videos no se almacenan en ningún
momento: cada fotograma se procesa para extraer las coordenadas de los puntos anatómicos y
la imagen se descarta, de modo que el conjunto de datos está compuesto únicamente por
valores numéricos a partir de los cuales no es posible reconstruir el rostro ni la
apariencia de una persona. La segunda es que todo el procesamiento ocurre en el propio
dispositivo, sin que ninguna imagen se transmita a servidores externos.

La participación de los colaboradores se documenta mediante una hoja de información y un
consentimiento informado, elaborados conforme a los principios éticos internacionales
aplicables a la investigación con seres humanos. El consentimiento distingue de forma separada la participación en el estudio, la
publicación del conjunto de datos bajo seudónimo y el uso de fotografías de las sesiones,
de manera que cada persona pueda autorizar unas y no otras.

La aplicación de estos marcos jurídicos y estándares técnicos asegura que el sistema
propuesto no solo responde a una necesidad tecnológica, sino que se alinea con las
obligaciones del Estado hondureño en materia de inclusión educativa, con los compromisos
internacionales adquiridos por el país y con las mejores prácticas de accesibilidad
digital.
